% 封面和书名页增强功能展示
% 演示 webtex-cn 对封面和书名页的丰富样式支持
\documentclass[极简]{ltc-guji}

% ========================================
% 示例 1: 基础封面（底色 + 文本框）
% ========================================
\begin{封面}[底色={244, 241, 225}]
  \文本框[
    floating=true,
    x=8cm,
    y=5cm,
    font-size=48pt,
    height=10
  ]{基础封面示例}

  \文本框[
    floating=true,
    x=10cm,
    y=12cm,
    font-size=20pt,
    height=3
  ]{简单底色}
\end{封面}

% ========================================
% 示例 2: 装饰封面（边框 + 内边距）
% ========================================
\begin{封面}[
  底色={245, 245, 220},
  边框=3pt solid,
  边框颜色={139, 69, 19},
  内边距=2cm
]
  \文本框[
    floating=true,
    x=8cm,
    y=7cm,
    font-size=42pt,
    height=8
  ]{装饰封面}

  \文本框[
    floating=true,
    x=10cm,
    y=13cm,
    font-size=18pt,
    height=4
  ]{带边框和内边距}
\end{封面}

% ========================================
% 示例 3: 基础书名页
% ========================================
\begin{书名页}
  \行[font-size=48pt, grid-height=52pt]{欽定四庫全書}
  \行[font-size=32pt, grid-height=36pt, align=center]{經部目錄}
  \行[font-size=24pt, grid-height=28pt, align=bottom]{卷一}
\end{书名页}

% ========================================
% 示例 4: 增强书名页（行间距 + 底色）
% ========================================
\begin{书名页}[
  行间距=20px,
  底色={250, 240, 230},
  内边距=3cm 2cm
]
  \行[font-size=48pt, grid-height=52pt, 字重=bold]{史記}
  \行[
    font-size=32pt,
    grid-height=36pt,
    align=center,
    颜色={139, 0, 0},
    字间距=0.2em
  ]{太史公書}
  \行[
    font-size=24pt,
    grid-height=28pt,
    align=bottom
  ]{卷一·本紀第一}
\end{书名页}

% ========================================
% 示例 5: 装饰线书名页
% ========================================
\begin{书名页}[行间距=15px, 内边距=2cm]
  \行[
    font-size=44pt,
    grid-height=48pt,
    装饰线=双边,
    内边距=0 1em
  ]{紅樓夢}

  \行[
    font-size=28pt,
    grid-height=32pt,
    align=center,
    装饰线=左,
    颜色={178, 34, 34}
  ]{甲戌本}

  \行[
    font-size=22pt,
    grid-height=26pt,
    align=bottom,
    字间距=0.15em
  ]{第一回}
\end{书名页}

% ========================================
% 示例 6: 自定义边框装饰
% ========================================
\begin{书名页}[底色={255, 250, 240}]
  \行[
    font-size=40pt,
    grid-height=44pt,
    边框=2px dashed,
    内边距=0.5em,
    字重=bold
  ]{詩經}

  \行[
    font-size=30pt,
    grid-height=34pt,
    align=center,
    颜色={70, 130, 180},
    margin=1em 0
  ]{國風·周南}

  \行[
    font-size=24pt,
    grid-height=28pt,
    align=bottom
  ]{關雎}
\end{书名页}

\begin{正文}
\chapter{正文示例}
这是正文内容，用于验证封面和书名页不会影响正文布局。
\end{正文}
